\section{Introducción}

El laboratorio solicita desplegar WordPress en una arquitectura de dos capas. La capa web ejecutará Apache, PHP y WordPress, mientras que una segunda máquina virtual alojará MySQL. La administración de ambos servidores se realizará desde una VM cliente mediante SSH y SFTP \cite{guia-lab5}.

El equipo eligió Proxmox VE como hipervisor Tipo 1. Esta decisión solo será válida para los puntos extra si Proxmox se instala directamente sobre hardware físico. La aclaración del docente excluye la virtualización anidada y prohíbe Terraform para este laboratorio.

\subsection{Estado de este documento}

Este PDF constituye un borrador estructural. Contiene decisiones, plantillas y criterios de prueba, pero no afirma que la infraestructura ya funcione. Las capturas, versiones efectivas, resultados, problemas y conclusiones finales se incorporarán después de ejecutar la implementación.

\pendiente{Completar el resumen de resultados cuando las pruebas T01 a T20 tengan evidencia real.}
